% Academic content with formatting based solely on the supplied Stanford guide.
% XeLaTeX; no URL commands, hyperlinks, tables, columns, or graphic elements.
\documentclass[10pt,a4paper]{article}
\usepackage[margin=1in]{geometry}
\usepackage{fontspec}
\setmainfont{texgyretermes}[Extension=.otf,UprightFont=*-regular,BoldFont=*-bold,ItalicFont=*-italic,BoldItalicFont=*-bolditalic,Ligatures=NoCommon]
\usepackage[protrusion=true,expansion=false]{microtype}
\usepackage{enumitem,titlesec,needspace}
\setlength{\parindent}{0pt}
\setlength{\parskip}{0pt}
\setlist[itemize]{leftmargin=12pt,itemsep=2pt,topsep=3pt,parsep=0pt,partopsep=0pt}
\titleformat{\section}{\normalfont\bfseries\fontsize{10.5}{12.2}\selectfont}{}{0pt}{}
\titlespacing*{\section}{0pt}{10pt}{5pt}
\pagestyle{plain}
\setlength{\emergencystretch}{2em}
\newcommand{\entry}[3]{\Needspace{3\baselineskip}\textbf{#1}\hfill #2\par #3\par\vspace{4pt}}
\newcommand{\paper}[4]{\begin{samepage}#1. #2. #3. #4\par\end{samepage}\vspace{6pt}}
\newcommand{\me}{\mbox{\textbf{Emmanuel G. Maminta}}}
\begin{document}
\fontsize{10.5}{12.2}\selectfont
\raggedright
\hyphenpenalty=10000
\exhyphenpenalty=10000
\special{pdf:docinfo << /Author (Emmanuel G. Maminta) /Title (Emmanuel G. Maminta - Curriculum Vitae) /Subject (Academic research and professional experience) >>}
{\bfseries\fontsize{17}{20}\selectfont EMMANUEL G. MAMINTA}\par
\vspace{4pt}
Metro Manila, Philippines\par
+63 945 113 2524\quad |\quad egmaminta@up.edu.ph\quad |\quad mnlmmnt@gmail.com

\section{EDUCATION}
\textbf{University of the Philippines Diliman}, Quezon City, Philippines\par
Ph.D. in Artificial Intelligence, Sep 2022–Present; expected completion 2027\par
Adviser: Prof. Rowel O. Atienza; Ubiquitous Computing Laboratory\par
M.Eng. in Artificial Intelligence, Sep 2022–Jul 2025\par
B.S. in Electrical Engineering, with honors, Aug 2018–Jul 2022

\section{RESEARCH INTERESTS}
Embodied AI, multimodal AI, and vision-language models for retail applications (RetailVision and Physical Retail AI); fine-grained grocery product retrieval and embodied agents in virtual retail environments. Additional areas of study: hand–object interaction, egocentric and exocentric computer vision, robotics, and world models.

\section{RESEARCH EXPERIENCE}
\textbf{Doctoral Researcher, Ubiquitous Computing Laboratory}\hfill Sep 2022–Present\par
University of the Philippines Diliman, Quezon City, Philippines
\begin{itemize}
  \item Co-authored a zero-shot evaluation of \textbf{190 open-source vision-language models} for grocery product retrieval, studying pre-training data, architecture, and input resolution (ICPR 2026).
  \item Led research on an agentic framework for embodied AI in virtual retail; co-authored Sari Sandbox, a simulation with over 250 interactive grocery items for benchmarking embodied agents.
  \item Investigated open-vocabulary object detection in vision-language models through research supported by Samsung Research \& Development Institute Philippines.
\end{itemize}

\section{PUBLICATIONS}
\paper{\me{} and Rowel O. Atienza}
{2026}
{A Transferable Variant Signal That Grocery Retrieval Cannot Reach}
{Asian Conference on Computer Vision (ACCV 2026), \textbf{Main Conference}. \textbf{Accepted for oral presentation}, Osaka, Japan, December 2026.}

\paper{\me{} and Rowel O. Atienza}
{2026}
{Late Interaction Ranking for Effective Closed-Catalog Grocery Product Retrieval}
{European Conference on Computer Vision (ECCV 2026), RetailVision Workshop. Presented in Malmö, Sweden, September 2026.}

\paper{\me{} and Rowel O. Atienza}
{2026}
{What Matters for Grocery Product Retrieval with Open Source Vision Language Models}
{International Conference on Pattern Recognition (ICPR 2026). Presented in Lyon, France, August 2026. \textbf{RRPR reproducibility badge}.}

\paper{Janika Deborah Gajo, Gerarld Paul Merales, Jerome Escarcha, Brenden Ashley Molina, Gian Nartea, \me, Juan Carlos Roldan, and Rowel O. Atienza}
{2025}
{Sari Sandbox: A Virtual Retail Store Environment for Embodied AI Agents}
{IEEE/CVF International Conference on Computer Vision (ICCV 2025), RetailVision Workshop. Presented in Honolulu, Hawai‘i, USA, October 2025.}

\section{HONORS AND AWARDS}
\entry{Most Active Contributor Award, Microsoft Research Asia}{Aug 2026}{Awarded during the AI Summer School at the National University of Singapore. Submitted the most contributions to the Global AI Value Challenge during the summer school.}
\entry{Best of ICCV 2025, Voxel51 feature}{Nov 2025}{Sari Sandbox was featured in Voxel51’s Best of ICCV 2025 research event.}

\newpage
\section{PROFESSIONAL EXPERIENCE}
\textbf{AI Developer, OpenText Philippines}\hfill Jan 2025–Present
\begin{itemize}
  \item Build enterprise AI products, agentic systems, and LLMOps/MLOps infrastructure; apply parameter-efficient fine-tuning (PEFT) and harness engineering.
  \item Developed an in-house small language model to automate proprietary coding tasks for teams in the Philippines and the United States.
  \item Engineered knowledge graph-guided retrieval-augmented generation (RAG) with a unified data model for internal applications.
\end{itemize}
\vspace{3pt}
\textbf{Data Scientist, Bank of the Philippine Islands}\hfill Nov 2023–Dec 2024
\begin{itemize}
  \item Engineered personalized product recommendations for high-net-worth clients and Jobs-to-Be-Done-based unsupervised clustering to identify customer segments.
\end{itemize}
\vspace{3pt}
\textbf{Software Engineer, Samsung Research \& Development Institute Philippines}\par
Apr–Oct 2023
\begin{itemize}
  \item Built a proprietary library for efficient, scalable deep learning experimentation.
\end{itemize}

\section{OPEN-SOURCE CONTRIBUTIONS}
\begin{itemize}
  \item \textbf{OpenMPR, Initiator.} Started an open-source project sharing recipes and evaluation resources for fine-grained grocery product retrieval.
  \item \textbf{stable-worldmodel, Contributor.} Contributed to a platform for reproducible world model research and evaluation.
  \item \textbf{GEPA-Lite, Creator.} Released a lightweight implementation of Genetic-Pareto prompt optimization under the MIT License.
  \item \textbf{GEPA, Contributor.} Added the AnyMaths adapter for prompt optimization on mathematical word problems and LLM reasoning.
\end{itemize}

\section{TEACHING AND RESEARCH COMMUNITY}
\entry{Cohere Labs Open Science Community, Member}{Sep 2026–Present}{Focus on multimodal AI and world models.}
\entry{Teaching contributions, University of the Philippines Diliman}{2022}{Co-taught EEE 111 (Introduction to Programming and Computation) and EEE 121 (Data Structures and Algorithms) with Prof. Jhoanna Rhodette I. Pedrasa.}

\section{CERTIFICATIONS AND TRAINING}
NVIDIA-Certified Professional: Agentic AI\par
NVIDIA-Certified Associate: Generative AI LLMs\par
Oxford Machine Learning Summer School, Jun–Jul 2024; enrolled in representation learning and generative AI courses.

\section{TECHNICAL SKILLS}
\textbf{Modeling:} Python, PyTorch, Transformers (Hugging Face), OpenCLIP, PEFT, OpenCV.\par
\textbf{Agent systems and evaluation:} DSPy, LangGraph, LlamaIndex, FastAPI, MLflow, Weights \& Biases, Langfuse.\par
\textbf{Inference and compute:} vLLM, NVIDIA Triton, ONNX, Slurm, NVIDIA Jetson, AWS, Azure, Google Cloud.
\end{document}
